\documentclass[10pt,a4paper]{amsart}
\usepackage[margin=19mm]{geometry}
\usepackage{amsmath,amssymb,mathtools}
\usepackage[scr=rsfs,cal=euler]{mathalfa}
\usepackage{xfrac}
\usepackage[unicode,hidelinks,bookmarks=false]{hyperref}
\newcommand{\ie}{i.e.\ }
\newcommand{\cf}{cf.\ }
\newcommand{\resp}{respectively\ }
\newcommand{\Subsection}{\subsection}
\providecommand{\nobreakdash}{\nobreak\hbox{-}}
\setlength{\emergencystretch}{2em}
\multlinegap0pt
\usepackage{bm}
\usepackage{bbm}
\usepackage{dsfont}
\usepackage{xspace}
\usepackage{cancel}
\usepackage{mathtools}
\usepackage{amsthm}
\makeatletter
\@ifundefined{equation}{\newtheorem{equation}{Equation}}{}
\@ifundefined{definition}{\newtheorem{definition}[equation]{Definition}}{}
\@ifundefined{lemma}{\newtheorem{lemma}[equation]{Lemma}}{}
\@ifundefined{remark}{\newtheorem{remark}[equation]{Remark}}{}
\@ifundefined{thm}{\newtheorem{thm}[equation]{Theorem}}{}
\makeatother
\newcommand{\R}{\mathbb{R}}
\newcommand{\calH}{\mathcal{H}}
\title[Restricted Projections to Hyperplanes in $\mathbb{R}^n$]{Restricted Projections to Hyperplanes in $\mathbb{R}^n$}
\author{Jiayin Liu}
\date{Source: arXiv:2604.14662v2 (CC BY 4.0)}
\begin{document}
\maketitle

\noindent\emph{Source and licence.} Restricted Projections to Hyperplanes in $\mathbb{R}^n$. Version arXiv:2604.14662v2 (CC BY 4.0), distributed under the Creative Commons Attribution 4.0 International licence, as recorded in the arXiv licence metadata for this exact version and shown on the abstract page https://arxiv.org/abs/2604.14662v2. Full rights notice: licenses/arxiv-2604.14662-CC-BY-4.0.txt.\par\smallskip

\noindent\textit{Editorial scope.} Counted unit: the Lemma whose source label is \texttt{cineprop} in the article (source0.tex lines 903--905, source label \texttt{cineprop}), delivered together with the article's own complete original proof of it (source0.tex lines 907--1010, one \texttt{proof} environment of 104 lines). No mathematics is written, repaired or re-derived: statement and proof are the article's own text line for line. Required context retained uncounted: planemain (lines 264--276); eix (lines 357--359); fw (lines 371--373); cine (lines 589--601); dual (lines 812--826); bdds (lines 866--868); II (lines 875--877); IIast (lines 879--881); basis (lines 885--900). Every definition, display and label of those lines is retained unchanged. Omitted from this package and not redistributed: 1--263, 277--356, 360--370, 374--588, 602--811, 827--865, 869--874, 878--878, 882--884, 901--902, 906--906, 1011--1863. Nothing inside a retained excerpt is omitted or altered: every retained body is delivered exactly as the article's lines are written, so no omission receipt of any kind accompanies this package. Citations: the retained excerpts carry no citation command at all, so no bibliography entry of the article is transcribed and no reference is redistributed. Reference adaptation: none was needed and none was made; every reference inside the delivered unit resolves inside it, so no unresolved reference or citation is produced. Printed slips kept literal: no printed word, symbol, hypothesis or number of the counted statement or of its proof was altered. Layout and class: the article's own class is \texttt{amsart} with packages including babel, fontenc, verbatim, palatino, amsmath, mathabx, amssymb, amsthm, amsfonts, graphicx, esint, color, mathtools, overpic; the wrapper is the lane's pinned \texttt{amsart} editorial wrapper at 10pt on A4, which restates the theorem-like environments the retained text needs and adds the article's own macro definitions that the retained text needs (\texttt{\textbackslash R}, \texttt{\textbackslash calH}), with extra wrapper packages (bm, bbm, dsfont, xspace, cancel, mathtools). The delivered numbering is the unit's own. Assets: no retained span contains a figure environment, a graphics-inclusion command, a PDF-page inclusion, a table or a TikZ picture; no image file is copied into this package, the package declares no asset dependency and no image file is redistributed with it. Licence: version arXiv:2604.14662v2 of this article is distributed under the Creative Commons Attribution 4.0 International licence, as recorded in the arXiv licence metadata for this exact version and shown on the abstract page https://arxiv.org/abs/2604.14662v2. A full rights notice is delivered at licenses/arxiv-2604.14662-CC-BY-4.0.txt. Characters: every retained excerpt is pure ASCII, so the delivered engine, XeLaTeX with its default Latin Modern text font, reports no missing character.
\begin{thm}\label{planemain}
For $n \ge 3$, let $Z$ be an analytic set in $\R^n$ with $\dim Z \le n-2$ and $\Sigma \subset S^{n-1}$ be an $(n-2)$-dimensional $C^2$ manifold having
\begin{itemize}
  \item non-vanishing geodesic curvature if $n=3$;
  \item sectional curvature $>1$ if $n \ge 4$.
\end{itemize}
Then, for any $s < \dim Z$, we have
  $$  \dim \{ x \in \Sigma : \dim [\pi_{T_{x}S^{n-1}}(Z)] < s  \} \le s.  $$
  In particular,
\begin{equation*}
   \dim [\pi_{T_{x}S^{n-1}}(Z)]  = \dim Z, \quad \calH^{n-2} \mbox{-\rm a.e.} \ x \in \Sigma.
\end{equation*}
\end{thm}

\begin{equation}\label{eix}
  e_i=e_i(x)=M_\Sigma^{-1}\nabla \Sigma(x)\cdot\partial_i, \quad \forall  i =1, \cdots, n-2
\end{equation}

\begin{equation}\label{fw}
  [0,1]^{n-2} \ni x \mapsto f_z(x):=(e_1(x)\cdot z, \cdots, e_{n-2}(x)\cdot z,  \nu(x)\cdot z) \in [0,1]^{n-1}
\end{equation}

\begin{definition}[Cinematic Map]\label{cine}
  Let $F \subset C^2(U,R^l)$ where $U \subset \mathbb{R}^k$ is a domain, $k,l \in \mathbb{N}$ with $l \ge k$. We say $F$ is a cinematic family of
maps, with cinematic constant $K$ and doubling constant $D$ if the following conditions
hold:
\begin{enumerate}
  \item $F$ is contained in a ball of diameter $K$ under the usual metric on $C^2(U,R^l)$;
  \item $F$ is a doubling metric space, with doubling constant at most $D$;
  \item For all $f, g \in F$, we have
  $$  \inf_{(x,\xi) \in U \times S^{k-1}} \{ |f(x) - g(x)| + |\nabla_\xi f(x) -\nabla_\xi g(x)| \} \ge K^{-1} \|f-g\|_{C^2(U)}  $$
  where $\nabla_\xi h= \nabla h \cdot \xi$ for $h=f,g$, noting that $\nabla h$ is a linear map from $\R^{k}$ to $\R^l$ and $\xi \in S^{k-1} \subset \R^{k}$ is considered as a direction of $\R^{k}$.
%  \item The family $\{\nabla_\xi\nabla_\zeta f \}_{f \in F, \xi,\zeta \in S^{k-1}} \subset C^0(U,\R^l)$ is equicontinuous.
\end{enumerate}
\end{definition}

\begin{lemma} \label{dual}
   Let $\Sigma \subset S^{n-1}$ be an $(n-2)$-dimensional $C^2$ manifold having
\begin{itemize}
  \item[(i)] non-vanishing geodesic curvature if $n=3$;
  \item[(ii)] sectional curvature $>1$ if $n \ge 4$.
\end{itemize}
  Then, $\Sigma^\ast$ also satisfies (i) and (ii) respectively.

  Moreover,
  \begin{enumerate}
  \item[(a)] if $\kappa_1$, $\cdots$, $\kappa_{n-2}$ are the geodesic curvature ($n=3$) or principal curvatures ($n\ge4$) at $x \in \Sigma$, then $\kappa_1^{-1}$, $\cdots$, $\kappa_{n-2}^{-1}$ are the geodesic curvature/principal curvatures at $\nu(x) \in \Sigma^\ast$;
  \item[(b)] for each $x \in \Sigma$, we have
  $$  T_{\nu(x)}\Sigma^\ast = T_{x}\Sigma  \mbox{ and } T_{\nu(x)}S^{n-1} = T_{\nu(x)}\Sigma^\ast \oplus {\rm span}\{x\}.  $$
  \end{enumerate}
\end{lemma}

\begin{equation}\label{bdds}
  \frac1{\bar\kappa}= \min_{x \in \Sigma^\ast,i =1,\cdots, n-2}\{\kappa_i^\ast(x)\} \mbox{ and } \frac1{\underline\kappa}= \max_{x \in \Sigma,i =1,\cdots, n-2}\{\kappa_i^\ast(x)\}.
\end{equation}

    \begin{equation}\label{II}
    (\nabla_\zeta \nu \cdot \zeta)(x) \le \bar \kappa|\zeta|^2, \quad \forall x \in \Sigma, \ \forall \zeta \in T_x\Sigma
\end{equation}

\begin{equation}\label{IIast}
    (\nabla_\zeta \nu^\ast \cdot \zeta)(x) \ge \frac{1}{\bar \kappa}|\zeta|^2, \quad \forall x \in \Sigma^\ast, \ \forall \zeta \in T_x\Sigma^\ast.
\end{equation}

\begin{remark}\label{basis}
Recall from \eqref{eix} that for each $x \in [0,1]^{n-2}$, $\{e_i(x)= \nabla\Sigma(x)\cdot\partial_i\}_{1 \le i \le n-2}$ is a basis of $T_x\Sigma$.
   Since $\Sigma:[0,1]^{n-2} \to \Sigma([0,1]^{n-2})$ is a diffeomorphism, for each $x \in [0,1]^{n-2}$, setting $w=\Sigma(x)$, we define
$$  \{\tilde e_i(w):= e_i(\Sigma^{-1}(w))= \nabla \Sigma(x) \cdot \partial_i\}_{1\le i \le n-2}. $$
Recalling Lemma \ref{dual} (b), we know that
\begin{equation*}
  {\rm span}  \{w=\nu^\ast, \tilde e_1, \cdots, \tilde e_{n-2}, \nu\}  = \mathbb{R}^{n}
\end{equation*}
where $\nu$ (resp. $\nu^\ast$) is the unit normal of $\Sigma \subset S^{n-1}$ (resp. $\Sigma^\ast \subset S^{n-1}$). In other words, $\{\nu^\ast, \tilde e_1, \cdots, \tilde e_{n-2}, \nu\}$ forms a basis of $\R^n$. Thus
we know for any $p \in \R^n$, $p$ can be uniquely written as
$$  p =  a_0 \nu^\ast + a_1 \tilde e_1 + \cdots + a_{n-2} \tilde e_{n-2} + a_{n-1} \nu
$$
where $a_i\in \R$ for $i=0, \cdots, n-1$. However, since the basis $\{\nu^\ast, \tilde e_1, \cdots, \tilde e_{n-2}, \nu\}$ may not be orthonormal, we have the following lower bound
$$  \sum_{i=0}^{n-1} a_i^2 \ge C_\Sigma |p|^2  $$
for some $0<C_\Sigma \le 1$ only depending on $\Sigma$.
\end{remark}

\begin{lemma} \label{cineprop}
  Let $Z$ be a subset in $B(o,1/2)\subset \mathbb{R}^{n}$ and $\Sigma : [0,1]^{n-2} \to \mathbb{S}^{n-1}$ be as in Theorem \ref{planemain}. Then  $F:=\{f_{z} \in C^2( [0,1]^{n-2},[0,1]^{n-1})\}_{z \in Z}$ defined in \eqref{fw} is a cinematic family with the same doubling constant as $\mathbb{R}^{n}$. In particular, $F$ is biLipschitz homeomorphic to $Z$.
\end{lemma}

\begin{proof}
First we show (1) in Definition \ref{cine}.
Since
\begin{equation}\label{cine1}
  \sup_{f_{z} \in F }\|f_{z}\|_{C^2([0,1]^{n-2})} \le \|\Sigma\|_{C^2([0,1]^{n-2})} \max_{z \in Z}|z| \lesssim 1
\end{equation}
 where the implicit constant only depends on $\Sigma$, (1) holds.

Fix $y,z \in Z$. We show (3) in Definition \ref{cine}. Similar as \eqref{cine1}, we have
\begin{equation}\label{dou1}
  \|f_y-f_{z}\|_{C^1([0,1]^{n-2})} \lesssim |y-z|.
\end{equation}
Write
$h:= f_y-f_{z}$. We prove (3) by showing
\begin{equation}\label{cine2}
  \inf_{x \in [0,1]^{n-2}, \xi \in S^{n-3}} \{ |h(x)| + |\nabla_\xi h(x)|  \} \gtrsim  |y-z|.
\end{equation}
We recall from Remark \ref{basis} that $w=\Sigma(x)$ and
$$  \{\tilde e_i(w):= e_i(\Sigma^{-1}(w))= \nabla \Sigma(x) \cdot \partial_i\}_{1\le i \le n-2} \mbox{ is a basis of $T_{w}\Sigma$}. $$
Define
$\tilde h  \in C^2(\Sigma)$ by
\begin{equation}\label{th}
  \tilde h(w):= ((y-z) \cdot \tilde e_1(w), \cdots,(y-z) \cdot \tilde e_{n-2}(w), \nu(w))  \quad w \in \Sigma.
\end{equation}
Thus $h=\tilde h \circ \Sigma$. For each $\xi \in S^{n-3}$, let $\xi_\ast:= \nabla \Sigma(x) \cdot \xi \in T_{\Sigma(x)} \Sigma$. We know for all $x \in [0,1]^{n-2}$ and $\xi \in S^{n-3}$,
\begin{equation}\label{rel3}
 \nabla_\xi h (x)= \langle\nabla h (x), \xi \rangle = \langle\nabla \tilde h (\Sigma(x)), \xi_\ast \rangle = \nabla_{\xi_\ast} \tilde h (w).
\end{equation}
and in particular,
\begin{equation} \label{chain}
  \nabla_{\partial_i} h (x) = \nabla_{\tilde e_i} \tilde h (w), \quad i =1, \cdots, n-2.
\end{equation}
Noting that $|\nabla \Sigma|$ is nowhere vanishing (since $\Sigma$ is a diffeomorphism), \eqref{rel3} implies that
\begin{equation}\label{rel4}
  |\nabla h (x)| \sim  |\nabla \tilde h (\Sigma(x))|, \quad x \in [0,1]^{n-2}
\end{equation}
where the implicit constant depends on $\max_{x \in [0,1]^{n-2}}\{|\nabla \Sigma(x)|\}$ and $\min_{x \in [0,1]^{n-2}}\{|\nabla \Sigma(x)|\}$.
Recalling again from Remark  \ref{basis} that
$y-z$ can be uniquely written as
$$  y-z =  a_0 \nu^\ast + a_1 \tilde e_1 + \cdots + a_{n-2} \tilde e_{n-2} + a_{n-1} \nu
$$
where $a_i\in \R$ for $i=0, \cdots, n-1$.
Moreover, recall that $\nabla_{\tilde e_i}\tilde e_j= \sum_{k=1}^{n-2}\Gamma_{i,j}^k \tilde e_k$ where $\Gamma_{i,j}^k$ are the Christoffel symbols determined by $\Sigma$. Let
\begin{equation}\label{M1}
  M_1:= \max_{w \in \Sigma}\max_{i,j,k=1, \cdots, n-2} |\Gamma_{i,j}^k(w)|.
\end{equation}
%Moreover, by Lemma \ref{dual}, we know $\Sigma^\ast$ has positive principal curvatures with the smallest one at least $\frac{1}{\bar \kappa}$, which implies that the second fundamental form
%\begin{equation}\label{II}
%   T\Sigma^\ast \ni \zeta \mapsto \nabla_\zeta \nu^\ast \cdot \zeta \ge \frac{1}{\bar \kappa}|\zeta|^2.
%\end{equation}
Let
$$  M_2 : = \frac{C_\Sigma}{16n M_1 \max\{1,\bar\kappa^2\}}$$
where $\bar\kappa$ is in \eqref{bdds} and $C_\Sigma$ is from Remark \ref{basis}.
We consider the following two cases separately.

\medskip
\emph{Case 1. $|a_i| \ge M_2|y-z|$ for some $i=1,\cdots, n-1$}. Then recalling \eqref{th},
we deduce that
$$ |h(x)| = |\tilde h (w)|  \ge   |(y-z) \cdot \tilde e_i| = |a_i| \ge M_2|y-z|  \sim |y-z|.$$
Thus \eqref{cine2} holds.

\medskip
\emph{Case 2. $|a_i| < M_2|y-z|$ for all $i=1,\cdots, n-1$}. Noting that $\Sigma_{i=0}^{n-1}a_i^2 \ge C_\Sigma |y-z|^2$, this implies that
\begin{equation}\label{a0}
  |a_0| \ge \frac12|y-z|.
\end{equation}
In this case, we show $|\nabla_\xi h(x)|\gtrsim |y-z|$. It suffices to check $\xi = \partial_i$ for $i=1, \cdots, n-2$, which, combining \eqref{chain} and \eqref{rel4}, further reduces to checking
$$   |\nabla_{\tilde e_i} \tilde h(w)|\gtrsim |y-z|, \quad  i= 1, \cdots,  n-2.   $$
By \eqref{M1}, we compute
\begin{align*}
  |(\nabla_{\tilde e_i}\tilde e_i)\cdot (y-z)| & \le \sum_{k=1}^{n-2} |\Gamma_{ii}^k \tilde e_k \cdot  (y-z)| \le \sum_{k=1}^{n-2} |\Gamma_{ii}^k a_k \tilde e_k| \le (n-2)M_1M_2|y-z|.
\end{align*}
Recalling $\tilde h$ in \eqref{th},
we have
\begin{equation*}
\nabla_{\tilde e_i} \tilde h (w)= (\nabla_{\tilde e_i}[(y-z) \cdot \tilde e_1], \cdots, \nabla_{\tilde e_i}[(y-z) \cdot \tilde e_{n-2}], \nabla_{\tilde e_i}[ (y-z) \cdot\nu]), \quad i=1, \cdots, n-2.
\end{equation*}
We compute
\begin{align*}
  |\nabla_{\tilde e_i} \tilde h(w)| & \ge |\nabla_{\tilde e_i}[(y-z) \cdot \tilde e_i]| =|[\nabla_{\tilde e_i}(y-z)]\cdot \tilde e_i + (\nabla_{\tilde e_i}\tilde e_i)\cdot (y-z)]|  \\
   & \ge  |\nabla_{\tilde e_i}[a_0\nu^\ast+[(y-z)-a_0\nu^\ast]]\cdot \tilde e_i| - |(\nabla_{\tilde e_i}\tilde e_i)\cdot(y-z)| \\
   & \ge |\nabla_{\tilde e_i}(a_0\nu^\ast)\cdot  \tilde e_i| - \sum_{j=1}^{n-2}|(\nabla_{\tilde e_i}a_je_j) \cdot \tilde e_i|- |(\nabla_{\tilde e_i}a_{n-1}\nu) \cdot \tilde e_i| - (n-2)M_1M_2|y-z|\\
   & \ge |a_0||\nabla_{\tilde e_i}\nu^\ast \cdot \tilde e_i| - \sum_{j=1}^{n-2}|a_j||(\nabla_{\tilde e_i}\tilde e_j) \cdot \tilde e_i| -|a_{n-1}||(\nabla_{\tilde e_i}\nu) \cdot \tilde e_i| - (n-2)M_1M_2|y-z|  \\
   & \ge \frac{1}{2 \bar \kappa}|y-z| - (n-2)M_1M_2 |y-z|- \bar \kappa M_2 |y-z| -(n-2)M_1M_2|y-z| \\
   & \ge \frac{1}{4 \bar \kappa}|y-z|
\end{align*}
where in the last two equations we use \eqref{a0}, \eqref{II} and \eqref{IIast}.
Thus \eqref{cine2} is true in this case. Combining two cases, we know (3) in Definition \ref{cine} holds.




\medskip

Finally, we show (2) in Definition \ref{cine}. Indeed, combining \eqref{dou1} and \eqref{cine2}, we deduce that
\begin{equation*}
 |y-z| \lesssim \inf_{x \in [0,1]^{n-2}, \xi \in S^{n-3}} \{ |h(x)| + |\nabla_\xi h(x)| \} \le  \|f_y-f_{z}\|_{C^1([0,1]^{n-1})} \lesssim |y-z|,
\end{equation*}
which implies that (2) in Definition \ref{cine} holds and
$F$ is a doubling space with the same doubling constant as $\mathbb{R}^{n}$. In particular, $F$ is biLipschitz homeomorphic to $Z$.


The proof is complete.
\end{proof}

\end{document}
